%! Independence, Basis, and Dimension — Unit 3, Lesson 3.4 — Slides (Beamer)
\documentclass[aspectratio=169, 11pt]{beamer}
\usepackage{linalg-beamer}   % provides \burgundyheader{} and \sectionlabel[color]{}
\newcommand{\vv}{\mathbf{v}}
\newcommand{\ww}{\mathbf{w}}
\newcommand{\xx}{\mathbf{x}}
\newcommand{\cc}{\mathbf{c}}
\newcommand{\bb}{\mathbf{b}}
\newcommand{\svec}{\mathbf{s}}
\newcommand{\zero}{\mathbf{0}}

\begin{document}

% TITLE SLIDE (CourseName is not defined in beamer, so the name is literal)
{
\setbeamercolor{background canvas}{bg=burgundy}
\begin{frame}[plain]
  \vspace{2.8cm}\hspace{0.55cm}%
  \begin{minipage}{0.9\textwidth}
    {\color{goldacc}\bfseries\small LESSON 3.4}\\[0.5cm]
    {\color{white}\bfseries\LARGE Independence, Basis, and Dimension}\\[1.0cm]
    {\color{blushmid}\small Linear Algebra~$\cdot$~Unit 3: The Four Fundamental Subspaces}
  \end{minipage}
\end{frame}
}

% HOOK
\begin{frame}[t, plain]
  \burgundyheader{How many levers does the arm \emph{really} have?}
  \sectionlabel{Hook}
  \begin{columns}[T]
    \begin{column}{0.58\textwidth}
      Each lever $\xx$ moves the tip to $A\xx$; the reachable set is $C(A)$.\\[6pt]
      The arm has \textbf{three} levers --- but lever~3 does exactly what levers~1 and~2 do together
      (column~3 $=$ column~1 $+$ column~2).\\[6pt]
      Only \emph{two} do anything new. The third is \textbf{redundant}.
    \end{column}
    \begin{column}{0.40\textwidth}
      \centering
      \begin{tikzpicture}[scale=0.55]
        \draw[->, thick, charcoal] (-0.3,0)--(3.2,0);
        \draw[->, thick, charcoal] (0,-0.3)--(0,3.2);
        \draw[->, very thick, burgundy] (0,0)--(2.4,0.8) node[right]{\scriptsize lever 1};
        \draw[->, very thick, royalblue] (0,0)--(0.8,2.4) node[above]{\scriptsize lever 2};
        \draw[->, very thick, goldacc, dashed] (0,0)--(3.2,3.2) node[right]{\scriptsize lever 3};
      \end{tikzpicture}
    \end{column}
  \end{columns}
  \vspace{2pt}
  \begin{block}{Three words for today}
    \emph{Independent} (no redundancy), \emph{basis} (just enough to reach everything), \emph{dimension}
    (how many independent directions).
  \end{block}
\end{frame}

% BIG IDEA: INDEPENDENCE = NULLSPACE
\begin{frame}[t, plain]
  \burgundyheader{Independence is a nullspace test}
  \sectionlabel{The big idea}
  Vectors are \textbf{independent} when the \emph{only} combination equal to $\zero$ is all-zero:
  \[ c_1\vv_1+\dots+c_k\vv_k=\zero \;\Rightarrow\; \text{every } c_i=0. \]
  Stack them as columns of $A$; a combination is $A\cc$, so:
  \begin{itemize}
    \setlength{\itemsep}{5pt}
    \item \textbf{Independent} $\iff N(A)=\{\zero\}$ (no free variables, $r=n$).
    \item \textbf{Dependent} $\iff$ a nonzero special solution $\svec$ --- which \emph{names} the redundancy.
  \end{itemize}
  \vspace{2pt}
  \begin{block}{Geometry}
    Two vectors are dependent when they share a line; three when they share a plane.
  \end{block}
\end{frame}

% READ IT OFF ONE REDUCTION
\begin{frame}[t, plain]
  \burgundyheader{Independence from one reduction}
  \sectionlabel[goldacc]{Test}
  \[
    A=\begin{bmatrix} 1 & 1 & 2 \\ 1 & 2 & 3 \end{bmatrix}
    \longrightarrow
    \begin{bmatrix} 1 & 0 & 1 \\ 0 & 1 & 1 \end{bmatrix}.
  \]
  \begin{itemize}
    \setlength{\itemsep}{4pt}
    \item Columns $1,2$ pivot; column $3$ is \textbf{free} $\Rightarrow$ a nonzero $\svec$ exists $\Rightarrow$ \textbf{dependent}.
    \item $\svec=\begin{bsmallmatrix} -1\\-1\\1 \end{bsmallmatrix}$ says $\mathbf{a}_3=\mathbf{a}_1+\mathbf{a}_2$: column~3 is redundant.
    \item Independence would need \emph{no} free column ($r=n$).
  \end{itemize}
\end{frame}

% BASIS
\begin{frame}[t, plain]
  \burgundyheader{Basis: independent \emph{and} spanning}
  \sectionlabel{Just enough, no waste}
  A \textbf{basis} spans the space \emph{and} is independent. From the same reduction:
  \begin{itemize}
    \setlength{\itemsep}{5pt}
    \item \textbf{Basis for $C(A)$:} the \emph{original} pivot columns --- $\begin{bsmallmatrix} 1\\1 \end{bsmallmatrix},\begin{bsmallmatrix} 1\\2 \end{bsmallmatrix}$ (drop the redundant column~3).
    \item \textbf{Basis for $N(A)$:} the special solutions --- $\svec=\begin{bsmallmatrix} -1\\-1\\1 \end{bsmallmatrix}$.
    \item \textbf{Standard basis} of $\mathbb{R}^n$: $\mathbf{e}_1,\dots,\mathbf{e}_n$.
  \end{itemize}
  \vspace{2pt}
  \begin{block}{Watch out}
    Use the pivot columns of $A$ itself --- \emph{not} the columns of the reduced matrix.
  \end{block}
\end{frame}

% DIMENSION
\begin{frame}[t, plain]
  \burgundyheader{Dimension --- independent directions}
  \sectionlabel[goldacc]{Count}
  Every basis of a space has the \emph{same} number of vectors: the \textbf{dimension}.
  \[
    \dim C(A)=r=2, \qquad \dim N(A)=n-r=1, \qquad r+(n-r)=n=3.
  \]
  \begin{itemize}
    \setlength{\itemsep}{4pt}
    \item Dimension counts the \emph{degrees of freedom} --- line $1$, plane $2$, all of $\mathbb{R}^3$ is $3$.
    \item More than $n$ vectors in $\mathbb{R}^n$ are always dependent (too many for the room).
    \item In $\mathbb{R}^n$, any $n$ independent vectors automatically span --- a basis.
  \end{itemize}
\end{frame}

% ROADMAP
\begin{frame}[t, plain]
  \burgundyheader{Where this goes}
  \sectionlabel{Connections}
  \textbf{Today:} independence ($N(A)=\{\zero\}$), basis (independent $+$ spans), and dimension
  ($\dim C(A)=r$, $\dim N(A)=n-r$) --- all from one reduction.\\[8pt]
  \begin{itemize}
    \setlength{\itemsep}{4pt}
    \item \textbf{3.5} Dimensions of the Four Subspaces: $C(A)$ and the row space have dimension $r$; $N(A)$ has $n-r$; the left nullspace $m-r$.
    \item \textbf{3.6} The Fundamental Theorem of Linear Algebra.
  \end{itemize}
  \vspace{4pt}
  \begin{block}{Keep in mind}
    One number $r$ controls it all: independent columns, the size of a basis, the dimension.
  \end{block}
\end{frame}

\end{document}
